\begin{textsample}{POS Dim 1 – 2022 – Score 32.00 – 2022/t003762.txt}  \label{ex:f1_pos_069}
Chris Anderson : Elon Musk, great to see you.
How are you?
Elon Musk : Good.
How are you?
CA : We ' re here at \textbf{the} Texas Gigafactory \textbf{the} day before this thing opens.
It ' s been pretty crazy out there.
Thank you so much for making time on a busy day.
I would love you to help us, kind of, cast our minds, I don ' t know, 10, 20, 30 years into \textbf{the} future.
And help us try to picture what it would take to build a future that ' s worth getting excited about. \textbf{The} last time you spoke at TED, you said that that was really just a big driver.
You know, you can talk about lots of other reasons to do \textbf{the} work you ' re doing, but fundamentally, you want to think about \textbf{the} future and not think that it sucks.
EM : Yeah, absolutely.
I think in general, you know, there ' s a lot of discussion of like, this problem or that problem.
And a lot of people are sad about \textbf{the} future and they ' re...
Pessimistic.
And I think... this is...
This is not great.
I mean, we really want to wake up in \textbf{the} morning and look forward to \textbf{the} future.
We want to be excited about what ' s going to happen.
And life cannot simply be about sort of, solving one miserable problem after another.
CA : So if you look forward 30 years, you know, \textbf{the} year 2050 has been labeled by scientists as this, kind of, almost like this doomsday deadline on climate.
There ' s a consensus of scientists, a large consensus of scientists, who believe that if we haven ' t completely eliminated \textbf{greenhouse} gases or offset them completely by 2050, effectively we ' re inviting climate catastrophe.
Do you believe there is a pathway to avoid that catastrophe?
And what would it look like?
EM : Yeah, so I am not one of \textbf{the} doomsday people, which may surprise you.
I actually think we ' re on a good path.
But at \textbf{the} same time, I want to caution against complacency.
So, so long as we are not complacent, as long as we have a high sense of urgency about moving towards a sustainable energy economy, then I think things will be fine.
So I can ' t emphasize that enough, as long as we push hard and are not complacent, \textbf{the} future is going to be great.
Don ' t worry about it.
I mean, worry about it, but if you worry about it, ironically, it will be a self-unfulfilling prophecy.
So, like, there are three elements to a sustainable energy future.
One is of sustainable energy generation, which is primarily wind and solar.
There ' s also hydro, geothermal, I ' m actually pro-nuclear.
I think nuclear is fine.
But it ' s going to be primarily solar and wind, as \textbf{the} primary generators of energy. \textbf{The} second part is you need batteries to store \textbf{the} solar and wind energy because \textbf{the} sun doesn ' t shine all \textbf{the} time, \textbf{the} wind doesn ' t blow all \textbf{the} time.
So it ' s a lot of stationary battery packs.
And then you need electric transport.
So electric cars, electric planes, boats.
And then ultimately, it ' s not really possible to make electric rockets, but you can make \textbf{the} propellant used in rockets using sustainable energy.
So ultimately, we can have a fully sustainable energy economy.
And it ' s those three things : solar / wind, stationary battery pack, electric vehicles.
So then what are \textbf{the} limiting factors on progress? \textbf{The} limiting factor really will be battery cell production.
So that ' s going to really be \textbf{the} fundamental rate driver.
And then whatever \textbf{the} slowest element of \textbf{the} whole lithium-ion battery cells supply chain, from mining and \textbf{the} many steps of refining to ultimately creating a battery cell and putting it into a pack, that will be \textbf{the} limiting factor on progress towards \textbf{sustainability}.
CA : All right, so we need to talk more about batteries, because \textbf{the} key thing that I want to understand, like, there seems to be a scaling issue here that is kind of amazing and alarming.
You have said that you have \textbf{calculated} that \textbf{the} amount of battery production that \textbf{the} world needs for \textbf{sustainability} is 300 terawatt hours of batteries.
That ' s \textbf{the} end goal?
EM : Very rough numbers, and I certainly would invite others to check our calculations because they may arrive at different conclusions.
But in order to transition, not just current electricity production, but also heating and transport, which roughly triples \textbf{the} amount of electricity that you need, it amounts to approximately 300 terawatt hours of installed capacity.
CA : So we need to give people a sense of how big a task that is.
I mean, here we are at \textbf{the} Gigafactory.
You know, this is one of \textbf{the} biggest buildings in \textbf{the} world.
What I ' ve read, and tell me if this is still right, is that \textbf{the} goal here is to eventually produce 100 gigawatt hours of batteries here a year eventually.
EM : We will probably do more than that, but yes, hopefully we get there within a couple of years.
CA : Right.
But I mean, that is one -- EM : 0. 1 terrawat hours.
CA : But that ' s still 1 / 100 of what ' s needed.
How much of \textbf{the} rest of that 100 is Tesla planning to take on let ' s say, between now and 2030, 2040, when we really need to see \textbf{the} scale up happen?
EM : I mean, these are just guesses.
So please, people shouldn ' t hold me to these things.
It ' s not like this is like some -- What tends to happen is I ' ll make some like, you know, best guess and then people, in five years, there ' ll be some jerk that writes an article :"Elon said this would happen, and it didn ' t happen.
He ' s a liar and a fool."It ' s very annoying when that happens.
So these are just guesses, this is a conversation.
CA : Right.
EM : I think Tesla probably ends up doing 10 percent of that.
Roughly.
CA : Let ' s say 2050 we have this amazing, you know, 100 percent sustainable electric grid made up of, you know, some mixture of \textbf{the} sustainable energy sources you talked about.
That same grid probably is offering \textbf{the} world really low-cost energy, isn ' t it, compared with now.
And I ' m curious about like, are people entitled to get a little bit excited about \textbf{the} possibilities of that world?
EM : People should be optimistic about \textbf{the} future.
Humanity will solve sustainable energy.
It will happen if we, you know, continue to push hard, \textbf{the} future is bright and good from an energy standpoint.
And then it will be possible to also use that energy to do \textbf{carbon} sequestration.
It takes a lot of energy to pull \textbf{carbon} out of \textbf{the} atmosphere because in putting it in \textbf{the} atmosphere it releases energy.
So now, you know, obviously in order to pull it out, you need to use a lot of energy.
But if you ' ve got a lot of sustainable energy from wind and solar, you can actually sequester \textbf{carbon}.
So you can reverse \textbf{the} CO2 parts per million of \textbf{the} atmosphere and oceans.
And also you can really have as much fresh water as you want.
Earth is mostly water.
We should call Earth"Water."It ' s 70 percent water by \textbf{surface} area.
Now most of that ' s seawater, but it ' s like we just happen to be on \textbf{the} bit that ' s land.
CA : And with energy, you can turn seawater into -- EM : Yes.
CA : Irrigating water or whatever water you need.
EM : At very low cost.
Things will be good.
CA : Things will be good.
And also, there ' s other benefits to this non-fossil fuel world where \textbf{the} air is cleaner -- EM : Yes, exactly.
Because, like, when you burn \textbf{fossil} fuels, there ' s all these side reactions and toxic gases of various kinds.
And sort of little particulates that are bad for your lungs.
Like, there ' s all sorts of bad things that are happening that will go away.
And \textbf{the} sky will be cleaner and quieter. \textbf{The} future ' s going to be good.
CA : I want us to switch now to think a bit about \textbf{artificial} intelligence.
But \textbf{the} segue there, you mentioned how annoying it is when people call you up for bad predictions in \textbf{the} past.
So I ' m possibly going to be annoying now, but I ' m curious about your timelines and how you predict and how come some things are so amazingly on \textbf{the} money and some aren ' t.
So when it comes to predicting sales of Tesla vehicles, for example, you ' ve kind of been amazing, I think in 2014 when Tesla had sold that year 60, 000 cars, you said,"2020, I think we will do half a million a year."EM : Yeah, we did almost exactly a half million.
CA : You did almost exactly half a million.
You were scoffed in 2014 because no one since Henry Ford, with \textbf{the} Model T, had come close to that kind of growth rate for cars.
You were scoffed, and you actually hit 500, 000 cars and then 510, 000 or whatever produced.
But five years ago, last time you came to TED, I asked you about full self-driving, and you said,"Yeah, this very year, I ' m confident that we will have a car going from LA to New York without any intervention."EM : Yeah, I don ' t want to blow your mind, but I ' m not always right.
CA : What ' s \textbf{the} difference between those two?
Why has full self-driving in particular been so hard to predict?
EM : I mean, \textbf{the} thing that really got me, and I think it ' s going to get a lot of other people, is that there are just so many false dawns with self-driving, where you think you ' ve got \textbf{the} problem, have a handle on \textbf{the} problem, and then it, no, turns out you just hit a ceiling.
Because if you were to plot \textbf{the} progress, \textbf{the} progress looks like a log curve.
So it ' s like a series of log curves.
So most people don ' t know what a log curve is, I suppose.
CA : Show \textbf{the} shape with your hands.
EM : It goes up you know, sort of a fairly straight way, and then it starts tailing off and you start getting diminishing returns.
And you ' re like, uh oh, it was trending up and now it ' s sort of, curving over and you start getting to these, what I call local maxima, where you don ' t realize basically how dumb you were.
And then it happens again.
And ultimately...
These things, you know, in retrospect, they seem obvious, but in order to solve full self-driving properly, you actually have to solve real-world AI.
Because what are \textbf{the} road networks designed to work with?
They ' re designed to work with a biological neural net, our brains, and with vision, our eyes.
And so in order to make it work with \textbf{computers}, you basically need to solve real-world AI and vision.
Because we need cameras and silicon neural nets in order to have self-driving work for a system that was designed for eyes and biological neural nets.
You know, I guess when you put it that way, it ' s sort of, like, quite obvious that \textbf{the} only way to solve full self-driving is to solve real world AI and sophisticated vision.
CA : What do you feel about \textbf{the} current architecture?
Do you think you have an architecture now where there is a chance for \textbf{the} logarithmic curve not to tail off any anytime soon?
EM : Well I mean, admittedly these may be infamous last words, but I actually am confident that we will solve it this year.
That we will exceed -- \textbf{The} \textbf{probability} of an accident, at what point do you exceed that of \textbf{the} average person?
I think we will exceed that this year.
CA : What are you seeing behind \textbf{the} scenes that gives you that confidence?
EM : We ' re almost at \textbf{the} point where we have a high-quality unified vector space.
In \textbf{the} beginning, we were trying to do this with image recognition on individual images.
But if you get one image out of a video, it ' s actually quite hard to see what ' s going on without ambiguity.
But if you look at a video segment of a few seconds of video, that ambiguity resolves.
So \textbf{the} first thing we had to do is tie all eight cameras together so they ' re synchronized, so that all \textbf{the} frames are looked at simultaneously and labeled simultaneously by one person, because we still need human labeling.
So at least they ' re not labeled at different times by different people in different ways.
So it ' s sort of a surround picture.
Then a very important part is to add \textbf{the} time \textbf{dimension}.
So that you ' re looking at surround video, and you ' re labeling surround video.
And this is actually quite difficult to do from a software standpoint.
We had to write our own labeling tools and then create auto labeling, create auto labeling software to amplify \textbf{the} efficiency of human labelers because it ' s quite hard to label.
In \textbf{the} beginning, it was taking several hours to label a 10-second video \textbf{clip}.
This is not scalable.
So basically what you have to have is you have to have surround video, and that surround video has to be primarily automatically labeled with humans just being editors and making slight corrections to \textbf{the} labeling of \textbf{the} video and then feeding back those corrections into \textbf{the} future auto labeler, so you get this flywheel eventually where \textbf{the} auto labeler is able to take in vast amounts of video and with high accuracy, automatically label \textbf{the} video for cars, lane lines, drive space.
CA : What you ' re saying is... \textbf{the} result of this is that you ' re effectively giving \textbf{the} car a 3D model of \textbf{the} actual \textbf{objects} that are all around it.
It knows what they are, and it knows how fast they are moving.
And \textbf{the} remaining task is to predict what \textbf{the} quirky behaviors are that, you know, that when a pedestrian is walking down \textbf{the} road with a smaller pedestrian, that maybe that smaller pedestrian might do something unpredictable or things like that.
You have to build into it before you can really call it safe.
EM : You basically need to have memory across time and space.
So what I mean by that is...
Memory can ' t be infinite, because it ' s using up a lot of \textbf{the} \textbf{computer} ' s RAM basically.
So you have to say how much are you going to try to remember?
It ' s very common for things to be occluded.
So if you talk about say, a pedestrian walking past a truck where you saw \textbf{the} pedestrian start on one side of \textbf{the} truck, then they ' re occluded by \textbf{the} truck.
You would know intuitively, OK, that pedestrian is going to pop out \textbf{the} other side, most likely.
CA : A \textbf{computer} doesn ' t know it.
EM : You need to slow down.
CA : A skeptic is going to say that every year for \textbf{the} last five years, you ' ve kind of said, well, no this is \textbf{the} year, we ' re confident that it will be there in a year or two or, you know, like it ' s always been about that far away.
But we ' ve got a new architecture now, you ' re seeing enough improvement behind \textbf{the} scenes to make you not certain, but pretty confident, that, by \textbf{the} end of this year, what in most, not in every city, and every circumstance but in many cities and circumstances, basically \textbf{the} car will be able to drive without interventions safer than a human.
EM : Yes.
I mean, \textbf{the} car currently drives me around Austin most of \textbf{the} time with no interventions.
So it ' s not like...
And we have over 100, 000 people in our full self-driving beta program.
So you can look at \textbf{the} videos that they post online.
CA : I do.
And some of them are great, and some of them are a little \textbf{terrifying}.
I mean, occasionally \textbf{the} car seems to veer off and scare \textbf{the} hell out of people.
EM : It ' s still a beta.
CA : But you ' re behind \textbf{the} scenes, looking at \textbf{the} data, you ' re seeing enough improvement to believe that a this-year timeline is real.
EM : Yes, that ' s what it seems like.
I mean, we could be here talking again in a year, like, well, another year went by, and it didn ' t happen.
But I think this is \textbf{the} year.
CA : And so in general, when people talk about Elon time, I mean it sounds like you can ' t just have a general rule that if you predict that something will be done in six months, actually what we should imagine is it ' s going to be a year or it ' s like two-x or three-x, it depends on \textbf{the} type of prediction.
Some things, I guess, things involving software, AI, whatever, are fundamentally harder to predict than others.
Is there an element that you actually deliberately make aggressive prediction timelines to drive people to be ambitious?
Without that, nothing gets done?
EM : Well, I generally believe, in terms of internal timelines, that we want to set \textbf{the} most aggressive timeline that we can.
Because there ' s sort of like a law of gaseous expansion where, for schedules, where whatever time you set, it ' s not going to be less than that.
It ' s very rare that it ' ll be less than that.
But as far as our predictions are concerned, what tends to happen in \textbf{the} media is that they will report all \textbf{the} wrong ones and ignore all \textbf{the} right ones.
Or, you know, when writing an article about me -- I ' ve had a long career in multiple industries.
If you list my sins, I sound like \textbf{the} worst person on Earth.
But if you put those against \textbf{the} things I ' ve done right, it makes much more sense, you know?
So essentially like, \textbf{the} longer you do anything, \textbf{the} more mistakes that you will make cumulatively.
Which, if you \textbf{sum} up those mistakes, will sound like I ' m \textbf{the} worst predictor ever.
But for example, for Tesla vehicle growth, I said I think we ' d do 50 percent, and we ' ve done 80 percent.
CA : Yes.
EM : But they don ' t mention that one.
So, I mean, I ' m not sure what my \textbf{exact} track record is on predictions.
They ' re more optimistic than pessimistic, but they ' re not all optimistic.
Some of them are exceeded probably more or later, but they do come true.
It ' s very rare that they do not come true.
It ' s sort of like, you know, if there ' s some radical technology prediction, \textbf{the} point is not that it was a few years late, but that it happened at all.
That ' s \textbf{the} more important part.
CA : So it feels like at some point in \textbf{the} last year, seeing \textbf{the} progress on understanding, \textbf{the} Tesla AI understanding \textbf{the} world around it, led to a kind of, an aha moment at Tesla.
Because you really surprised people recently when you said probably \textbf{the} most important product development going on at Tesla this year is this robot, Optimus.
EM : Yes.
CA : Many companies out there have tried to put out these robots, they ' ve been working on them for years.
And so far no one has really cracked it.
There ' s no mass adoption robot in people ' s homes.
There are some in manufacturing, but I would say, no one ' s kind of, really cracked it.
Is it something that happened in \textbf{the} development of full self-driving that gave you \textbf{the} confidence to say,"You know what, we could do something special here."EM : Yeah, exactly.
So, you know, it took me a while to sort of realize that in order to solve self-driving, you really needed to solve real-world AI.
And at \textbf{the} point of which you solve real-world AI for a car, which is really a robot on four wheels, you can then generalize that to a robot on legs as well. \textbf{The} two hard parts I think -- like obviously companies like Boston Dynamics have shown that it ' s possible to make quite compelling, sometimes alarming robots.
CA : Right.
EM : You know, so from a sensors and actuators standpoint, it ' s certainly been demonstrated by many that it ' s possible to make a humanoid robot. \textbf{The} things that are currently missing are enough intelligence for \textbf{the} robot to navigate \textbf{the} real world and do useful things without being explicitly instructed.
So \textbf{the} missing things are basically real-world intelligence and scaling up manufacturing.
Those are two things that Tesla is very good at.
And so then we basically just need to design \textbf{the} specialized actuators and sensors that are needed for humanoid robot.
People have no idea, this is going to be bigger than \textbf{the} car.
CA : So let ' s dig into exactly that.
I mean, in one way, it ' s actually an easier problem than full self-driving because instead of an \textbf{object} going along at 60 miles an hour, which if it gets it wrong, someone will die.
This is an \textbf{object} that ' s engineered to only go at what, three or four or five miles an hour.
And so a mistake, there aren ' t lives at stake.
There might be embarrassment at stake.
EM : So long as \textbf{the} AI doesn ' t take it over and murder us in our sleep or something.
CA : Right.
So talk about -- I think \textbf{the} first applications you ' ve mentioned are probably going to be manufacturing, but eventually \textbf{the} vision is to have these available for people at home.
If you had a robot that really understood \textbf{the} 3D architecture of your house and knew where every \textbf{object} in that house was or was supposed to be, and could recognize all those \textbf{objects}, I mean, that ' s kind of amazing, isn ' t it?
Like \textbf{the} kind of thing that you could ask a robot to do would be what?
Like, tidy up?
EM : Yeah, absolutely.
Make dinner, I guess, mow \textbf{the} lawn.
CA : Take a cup of tea to grandma and show her family pictures.
EM : Exactly.
Take care of my grandmother and make sure -- CA : It could obviously recognize everyone in \textbf{the} home.
It could play catch with your kids.
EM : Yes.
I mean, obviously, we need to be careful this doesn ' t become a dystopian situation.
I think one of \textbf{the} things that ' s going to be important is to have a localized ROM chip on \textbf{the} robot that cannot be updated over \textbf{the} air.
Where if you, for example, were to say,"Stop, stop, stop,"if anyone said that, then \textbf{the} robot would stop, you know, type of thing.
And that ' s not updatable remotely.
I think it ' s going to be important to have safety features like that.
CA : Yeah, that sounds wise.
EM : And I do think there should be a regulatory agency for AI.
I ' ve said that for many years.
I don ' t love being regulated, but I think this is an important thing for public safety.
CA : Let ' s come back to that.
But I don ' t think many people have really sort of taken seriously \textbf{the} notion of, you know, a robot at home.
I mean, at \textbf{the} start of \textbf{the} computing revolution, Bill Gates said there ' s going to be a \textbf{computer} in every home.
And people at \textbf{the} time said, yeah, whatever, who would even want that.
Do you think there will be basically like in, say, 2050 or whatever, like a robot in most homes, is what there will be, and people will love them and count on them?
You ' ll have your own butler basically.
EM : Yeah, you ' ll have your sort of buddy robot probably, yeah.
CA : I mean, how much of a buddy?
How many applications have you thought, you know, can you have a romantic partner, a sex partner?
EM : It ' s probably inevitable.
I mean, I did promise \textbf{the} internet that I ' d make catgirls.
We could make a robot catgirl.
CA : Be careful what you promise \textbf{the} internet.
EM : So, yeah, I guess it ' ll be whatever people want really, you know.
CA : What sort of timeline should we be thinking about of \textbf{the} first models that are actually made and sold?
EM : Well, you know, \textbf{the} first units that we intend to make are for jobs that are dangerous, boring, repetitive, and things that people don ' t want to do.
And, you know, I think we ' ll have like an interesting prototype sometime this year.
We might have something useful next year, but I think quite likely within at least two years.
And then we ' ll see rapid growth year over year of \textbf{the} usefulness of \textbf{the} humanoid robots and decrease in cost and scaling up production.
CA : Initially just selling to businesses, or when do you picture you ' ll start selling them where you can buy your parents one for Christmas or something?
EM : I ' d say in less than ten years.
CA : Help me on \textbf{the} economics of this.
So what do you picture \textbf{the} cost of one of these being?
EM : Well, I think \textbf{the} cost is actually not going to be crazy high.
Like less than a car.
Initially, things will be expensive because it ' ll be a new technology at low production volume. \textbf{The} complexity and cost of a car is greater than that of a humanoid robot.
So I would expect that it ' s going to be less than a car, or at least equivalent to a cheap car.
CA : So even if it starts at 50k, within a few years, it ' s down to 20k or lower or whatever.
And maybe for home they ' ll get much cheaper still.
But think about \textbf{the} economics of this.
If you can replace a \$30, 000, \$40, 000-a-year worker, which you have to pay every year, with a one-time payment of \$25, 000 for a robot that can work longer hours, a pretty rapid replacement of certain types of jobs.
How worried should \textbf{the} world be about that?
EM : I wouldn ' t worry about \textbf{the} sort of, putting people out of a job thing.
I think we ' re actually going to have, and already do have, a massive shortage of labor.
So I think we will have...
Not people out of work, but actually still a shortage labor even in \textbf{the} future.
But this really will be a world of abundance.
Any goods and services will be available to anyone who wants them.
It ' ll be so cheap to have goods and services, it will be ridiculous.
CA : I ' m presuming it should be possible to imagine a bunch of goods and services that can ' t profitably be made now but could be made in that world, courtesy of legions of robots.
EM : Yeah.
It will be a world of abundance. \textbf{The} only scarcity that will exist in \textbf{the} future is that which we decide to create ourselves as humans.
CA : OK.
So AI is allowing us to imagine a differently powered economy that will create this abundance.
What are you most worried about going wrong?
EM : Well, like I said, AI and robotics will bring out what might be termed \textbf{the} age of abundance.
Other people have used this word, and that this is my prediction : it will be an age of abundance for everyone.
But I guess there ' s... \textbf{The} dangers would be \textbf{the} \textbf{artificial} general intelligence or digital superintelligence decouples from a collective human will and goes in \textbf{the} direction that for some reason we don ' t like.
Whatever direction it might go.
You know, that ' s sort of \textbf{the} idea behind Neuralink, is to try to more tightly couple collective human world to digital superintelligence.
And also along \textbf{the} way solve a lot of brain injuries and spinal injuries and that kind of thing.
So even if it doesn ' t succeed in \textbf{the} greater goal, I think it will succeed in \textbf{the} goal of alleviating brain and spine damage.
CA : So \textbf{the} spirit there is that if we ' re going to make these AIs that are so vastly intelligent, we ought to be wired directly to them so that we ourselves can have those superpowers more directly.
But that doesn ' t seem to avoid \textbf{the} risk that those superpowers might... turn ugly in unintended ways.
EM : I think it ' s a risk, I agree.
I ' m not saying that I have some certain answer to that risk.
I ' m just saying like maybe one of \textbf{the} things that would be good for ensuring that \textbf{the} future is one that we want is to more tightly couple \textbf{the} collective human world to digital intelligence. \textbf{The} issue that we face here is that we are already a cyborg, if you think about it. \textbf{The} \textbf{computers} are an extension of ourselves.
And when we die, we have, like, a digital ghost.
You know, all of our text messages and social media, emails.
And it ' s quite eerie actually, when someone dies but everything online is still there.
But you say like, what ' s \textbf{the} limitation?
What is it that inhibits a human-machine symbiosis?
It ' s \textbf{the} data rate.
When you communicate, especially with a phone, you ' re moving your thumbs very slowly.
So you ' re like moving your two little meat sticks at a rate that ' s maybe 10 bits per second, optimistically, 100 bits per second.
And \textbf{computers} are communicating at \textbf{the} gigabyte level and beyond.
CA : Have you seen evidence that \textbf{the} technology is actually working, that you ' ve got a richer, sort of, higher bandwidth connection, if you like, between like external electronics and a brain than has been possible before?
EM : Yeah.
I mean, \textbf{the} fundamental principles of reading neurons, sort of doing read-write on neurons with tiny electrodes, have been demonstrated for decades.
So it ' s not like \textbf{the} concept is new. \textbf{The} problem is that there is no product that works well that you can go and buy.
So it ' s all sort of, in research labs.
And it ' s like some cords sticking out of your head.
And it ' s quite gruesome, and it ' s really...
There ' s no good product that actually does a good job and is high-bandwidth and safe and something actually that you could buy and would want to buy.
But \textbf{the} way to think of \textbf{the} Neuralink \textbf{device} is kind of like a Fitbit or an Apple Watch.
That ' s where we take out sort of a small section of skull about \textbf{the} size of a quarter, replace that with what, in many ways really is very much like a Fitbit, Apple Watch or some kind of smart watch thing.
But with tiny, tiny wires, very, very tiny wires.
Wires so tiny, it ' s hard to even see them.
And it ' s very important to have very tiny wires so that when they ' re implanted, they don ' t damage \textbf{the} brain.
CA : How far are you from putting these into humans?
EM : Well, we have put in our FDA application to aspirationally do \textbf{the} first human implant this year.
CA : \textbf{The} first uses will be for neurological injuries of different kinds.
But rolling \textbf{the} clock forward and imagining when people are actually using these for their own enhancement, let ' s say, and for \textbf{the} enhancement of \textbf{the} world, how clear are you in your mind as to what it will feel like to have one of these inside your head?
EM : Well, I do want to emphasize we ' re at an early stage.
And so it really will be many years before we have anything approximating a high-bandwidth neural interface that allows for AI-human symbiosis.
For many years, we will just be solving brain injuries and spinal injuries.
For probably a decade.
This is not something that will suddenly one day it will have this incredible sort of whole brain interface.
It ' s going to be, like I said, at least a decade of really just solving brain injuries and spinal injuries.
And really, I think you can solve a very wide range of brain injuries, including severe depression, morbid obesity, sleep, potentially schizophrenia, like, a lot of things that cause great stress to people.
Restoring memory in older people.
CA : If you can pull that off, that ' s \textbf{the} app I will sign up for.
EM : Absolutely.
CA : Please hurry.
EM : I mean, \textbf{the} emails that we get at Neuralink are heartbreaking.
I mean, they ' ll send us just tragic, you know, where someone was sort of, in \textbf{the} prime of life and they had an accident on a motorcycle and someone who ' s 25, you know, can ' t even feed themselves.
And this is something we could fix.
CA : But you have said that AI is one of \textbf{the} things you ' re most worried about and that Neuralink may be one of \textbf{the} ways where we can keep abreast of it.
EM : Yeah, there ' s \textbf{the} short-term thing, which I think is helpful on an individual human level with injuries.
And then \textbf{the} long-term thing is an attempt to address \textbf{the} civilizational risk of AI by bringing digital intelligence and biological intelligence closer together.
I mean, if you think of how \textbf{the} brain works today, there are really two layers to \textbf{the} brain.
There ' s \textbf{the} limbic system and \textbf{the} cortex.
You ' ve got \textbf{the} kind of, animal brain where -- it ' s kind of like \textbf{the} fun part, really.
CA : It ' s where most of Twitter operates, by \textbf{the} way.
EM : I think Tim Urban said, we ' re like somebody, you know, stuck a \textbf{computer} on a monkey.
You know, so we ' re like, if you gave a monkey a \textbf{computer}, that ' s our cortex.
But we still have a lot of monkey instincts.
Which we then try to rationalize as, no, it ' s not a monkey instinct.
It ' s something more important than that.
But it ' s often just really a monkey instinct.
We ' re just monkeys with a \textbf{computer} stuck in our brain.
But even though \textbf{the} cortex is sort of \textbf{the} smart, or \textbf{the} intelligent part of \textbf{the} brain, \textbf{the} thinking part of \textbf{the} brain, I ' ve not yet met anyone who wants to delete their limbic system or their cortex.
They ' re quite happy having both.
Everyone wants both parts of their brain.
And people really want their phones and their \textbf{computers}, which are really \textbf{the} tertiary, \textbf{the} third part of your intelligence.
It ' s just that it ' s...
Like \textbf{the} bandwidth, \textbf{the} rate of communication with that tertiary layer is slow.
And it ' s just a very tiny straw to this tertiary layer.
And we want to make that tiny straw a big \textbf{highway}.
And I ' m definitely not saying that this is going to solve everything.
Or this is you know, it ' s \textbf{the} only thing -- it ' s something that might be helpful.
And worst-case scenario, I think we solve some important brain injury, spinal injury issues, and that ' s still a great outcome.
CA : Best-case scenario, we may discover new human possibility, telepathy, you ' ve spoken of, in a way, a connection with a loved one, you know, full memory and much faster thought processing maybe.
All these things.
It ' s very cool.
If AI were to take down Earth, we need a plan B.
Let ' s shift our attention to space.
We spoke last time at TED about reusability, and you had just demonstrated that spectacularly for \textbf{the} first time.
Since then, you ' ve gone on to build this monster rocket, Starship, which kind of changes \textbf{the} rules of \textbf{the} game in spectacular ways.
Tell us about Starship.
EM : Starship is extremely fundamental.
So \textbf{the} holy grail of rocketry or space transport is full and rapid reusability.
This has never been achieved. \textbf{The} closest that anything has come is our Falcon 9 rocket, where we are able to recover \textbf{the} first stage, \textbf{the} boost stage, which is probably about 60 percent of \textbf{the} cost of \textbf{the} vehicle of \textbf{the} whole launch, maybe 70 percent.
And we ' ve now done that over a hundred times.
So with Starship, we will be recovering \textbf{the} entire thing.
Or at least that ' s \textbf{the} goal.
CA : Right.
EM : And moreover, recovering it in such a way that it can be immediately re-flown.
Whereas with Falcon 9, we still need to do some amount of refurbishment to \textbf{the} booster and to \textbf{the} fairing nose cone.
But with Starship, \textbf{the} design goal is immediate re-flight.
So you just refill propellants and go again.
And this is gigantic.
Just as it would be in any other mode of transport.
CA : And \textbf{the} main design is to basically take 100 plus people at a time, plus a bunch of things that they need, to Mars.
So, first of all, talk about that piece.
What is your latest timeline?
One, for \textbf{the} first time, a Starship goes to Mars, presumably without people, but just equipment.
Two, with people.
Three, there ' s sort of, OK, 100 people at a time, let ' s go.
EM : Sure.
And just to put \textbf{the} cost thing into perspective, \textbf{the} expected cost of Starship, putting 100 tons into orbit, is significantly less than what it would have cost or what it did cost to put our tiny Falcon 1 rocket into orbit.
Just as \textbf{the} cost of flying a 747 around \textbf{the} world is less than \textbf{the} cost of a small \textbf{airplane}.
You know, a small \textbf{airplane} that was thrown away.
So it ' s really pretty mind-boggling that \textbf{the} giant thing costs less, way less than \textbf{the} small thing.
So it doesn ' t use exotic propellants or things that are difficult to obtain on Mars.
It uses methane as fuel, and it ' s primarily oxygen, roughly 77-78 percent oxygen by weight.
And Mars has a CO2 atmosphere and has water ice, which is CO2 plus H2O, so you can make CH4, methane, and O2, oxygen, on Mars.
CA : Presumably, one of \textbf{the} first tasks on Mars will be to create a fuel plant that can create \textbf{the} fuel for \textbf{the} return trips of many Starships.
EM : Yes.
And actually, it ' s mostly going to be oxygen plants, because it ' s 78 percent oxygen, 22 percent fuel.
But \textbf{the} fuel is a simple fuel that is easy to create on Mars.
And in many other parts of \textbf{the} solar system.
So basically...
And it ' s all propulsive landing, no parachutes, nothing thrown away.
It has a heat shield that ' s capable of entering on Earth or Mars.
We can even potentially go to Venus. but you don ' t want to go there.
Venus is hell, almost literally.
But you could...
It ' s a generalized method of transport to anywhere in \textbf{the} solar system, because \textbf{the} point at which you have propellant depo on Mars, you can then travel to \textbf{the} \textbf{asteroid} belt and to \textbf{the} moons of Jupiter and Saturn and ultimately anywhere in \textbf{the} solar system.
CA : But your main focus and SpaceX ' s main focus is still Mars.
That is \textbf{the} mission.
That is where most of \textbf{the} effort will go?
Or are you actually imagining a much broader \textbf{array} of uses even in \textbf{the} coming, you know, \textbf{the} first decade or so of uses of this.
Where we could go, for example, to other places in \textbf{the} solar system to explore, perhaps NASA wants to use \textbf{the} rocket for that reason.
EM : Yeah, NASA is planning to use a Starship to return to \textbf{the} moon, to return people to \textbf{the} moon.
And so we ' re very honored that NASA has chosen us to do this.
But I ' m saying it is a generalized -- it ' s a general solution to getting anywhere in \textbf{the} greater solar system.
It ' s not suitable for going to another star system, but it is a general solution for transport anywhere in \textbf{the} solar system.
CA : Before it can do any of that, it ' s got to demonstrate it can get into orbit, you know, around Earth.
What ' s your latest advice on \textbf{the} timeline for that?
EM : It ' s looking promising for us to have an orbital launch attempt in a few months.
So we ' re actually integrating -- will be integrating \textbf{the} engines into \textbf{the} booster for \textbf{the} first orbital flight starting in about a week or two.
And \textbf{the} launch complex itself is ready to go.
So assuming we get regulatory approval, I think we could have an orbital launch attempt within a few months.
CA : And a radical new technology like this presumably there is real risk on those early attempts.
EM : Oh, 100 percent, yeah. \textbf{The} joke I make all \textbf{the} time is that excitement is guaranteed.
Success is not guaranteed, but excitement certainly is.
CA : But \textbf{the} last I saw on your timeline, you ' ve slightly put back \textbf{the} expected date to put \textbf{the} first human on Mars till 2029, I want to say?
EM : Yeah, I mean, so let ' s see.
I mean, we have built a production system for Starship, so we ' re making a lot of ships and boosters.
CA : How many are you planning to make actually?
EM : Well, we ' re currently expecting to make a booster and a ship roughly every, well, initially, roughly every couple of months, and then hopefully by \textbf{the} end of this year, one every month.
So it ' s giant rockets, and a lot of them.
Just talking in terms of rough orders of magnitude, in order to create a self-sustaining city on Mars, I think you will need something on \textbf{the} order of a thousand ships.
And we just need a Helen of Sparta, I guess, on Mars.
CA : This is not in most people ' s heads, Elon.
EM : \textbf{The} \textbf{planet} that launched 1, 000 ships.
CA : That ' s nice.
But this is not in most people ' s heads, this picture that you have in your mind.
There ' s basically a two-year window, you can only really fly to Mars conveniently every two years.
You were picturing that during \textbf{the} 2030s, every couple of years, something like 1, 000 Starships take off, each containing 100 or more people.
That picture is just completely mind-blowing to me.
That sense of this armada of humans going to -- EM : It ' ll be like"Battlestar Galactica,"\textbf{the} fleet departs.
CA : And you think that it can basically be funded by people spending maybe a couple hundred grand on a \textbf{ticket} to Mars?
Is that price about where it has been?
EM : Well, I think if you say like, what ' s required in order to get enough people and enough cargo to Mars to build a self-sustaining city.
And it ' s where you have an intersection of sets of people who want to go, because I think only a small percentage of humanity will want to go, and can afford to go or get sponsorship in some manner.
That intersection of sets, I think, needs to be a million people or something like that.
And so it ' s what can a million people afford, or get sponsorship for, because I think governments will also pay for it, and people can take out loans.
But I think at \textbf{the} point at which you say, OK, like, if moving to Mars costs are, for argument ' s sake, \$100, 000, then I think you know, almost anyone can work and save up and eventually have \$100, 000 and be able to go to Mars if they want.
We want to make it available to anyone who wants to go.
It ' s very important to emphasize that Mars, especially in \textbf{the} beginning, will not be luxurious.
It will be dangerous, cramped, difficult, hard work.
It ' s kind of like that Shackleton ad for going to \textbf{the} Antarctic, which I think is actually not real, but it sounds real and it ' s cool.
It ' s sort of like, \textbf{the} sales pitch for going to Mars is,"It ' s dangerous, it ' s cramped.
You might not make it back.
It ' s difficult, it ' s hard work."That ' s \textbf{the} sales pitch.
CA : Right.
But you will make history.
EM : But it ' ll be glorious.
CA : So on that kind of launch rate you ' re talking about over two decades, you could get your million people to Mars, essentially.
Whose city is it?
Is it NASA ' s city, is it SpaceX ' s city?
EM : It ' s \textbf{the} people of Mars ' city. \textbf{The} reason for this, I mean, I feel like why do this thing?
I think this is important for maximizing \textbf{the} probable lifespan of humanity or consciousness.
Human civilization could come to an end for external reasons, like a giant meteor or super volcanoes or extreme climate change.
Or World War III, or you know, any one of a number of reasons.
But \textbf{the} probable life span of civilizational consciousness as we know it, which we should really view as this very delicate thing, like a small candle in a vast darkness.
That is what appears to be \textbf{the} case.
We ' re in this vast darkness of space, and there ' s this little candle of consciousness that ' s only really come about after 4. 5 billion years, and it could just go out.
CA : I think that ' s powerful, and I think a lot of people will be inspired by that vision.
And \textbf{the} reason you need \textbf{the} million people is because there has to be enough people there to do everything that you need to survive.
EM : Really, like \textbf{the} critical threshold is if \textbf{the} ships from Earth stop coming for any reason, does \textbf{the} Mars City die out or not?
And so we have to -- You know, people talk about like, \textbf{the} sort of, \textbf{the} great filters, \textbf{the} things that perhaps, you know, we talk about \textbf{the} Fermi paradox, and where are \textbf{the} aliens?
Well maybe there are these various great filters that \textbf{the} aliens didn ' t pass, and so they eventually just ceased to exist.
And one of \textbf{the} great filters is becoming a multi-planet species.
So we want to pass that filter.
And I ' ll be long-dead before this is, you know, a real thing, before it happens.
But I ' d like to at least see us make great progress in this direction.
CA : Given how tortured \textbf{the} Earth is right now, how much we ' re beating each other up, shouldn ' t there be discussions going on with everyone who is dreaming about Mars to try to say, we ' ve got a once in a civilization ' s chance to make some new rules here?
Should someone be trying to lead those discussions to figure out what it means for this to be \textbf{the} people of Mars ' City?
EM : Well, I think ultimately this will be up to \textbf{the} people of Mars to decide how they want to rethink society.
Yeah there ' s certainly risk there.
And hopefully \textbf{the} people of Mars will be more enlightened and will not fight amongst each other too much.
I mean, I have some recommendations, which people of Mars may choose to listen to or not.
I would advocate for more of a direct democracy, not a representative democracy, and laws that are short enough for people to understand.
Where it is harder to create laws than to get rid of them.
CA : Coming back a bit nearer term, I ' d love you to just talk a bit about some of \textbf{the} other possibility space that Starship seems to have created.
So given -- Suddenly we ' ve got this ability to move 100 tons-plus into orbit.
So we ' ve just launched \textbf{the} James Webb telescope, which is an incredible thing.
It ' s unbelievable.
EM : Exquisite piece of technology.
CA : Exquisite piece of technology.
But people spent two years trying to figure out how to fold up this thing.
It ' s a three-ton telescope.
EM : We can make it a lot easier if you ' ve got more volume and mass.
CA : But let ' s ask a different question.
Which is, how much more powerful a telescope could someone design based on using Starship, for example?
EM : I mean, roughly, I ' d say it ' s probably an order of magnitude more resolution.
If you ' ve got 100 tons and a thousand cubic meters volume, which is roughly what we have.
CA : And what about other exploration through \textbf{the} solar system?
I mean, I ' m you know -- EM : Europa is a big question mark.
CA : Right, so there ' s an ocean there.
And what you really want to do is to drop a submarine into that ocean.
EM : Maybe there ' s like, some squid civilization, cephalopod civilization under \textbf{the} ice of Europa.
That would be pretty interesting.
CA : I mean, Elon, if you could take a submarine to Europa and we see pictures of this thing being devoured by a squid, that would honestly be \textbf{the} happiest moment of my life.
EM : Pretty wild, yeah.
CA : What other possibilities are out there?
Like, it feels like if you ' re going to create a thousand of these things, they can only fly to Mars every two years.
What are they doing \textbf{the} rest of \textbf{the} time?
It feels like there ' s this \textbf{explosion} of possibility that I don ' t think people are really thinking about.
EM : I don ' t know, we ' ve certainly got a long way to go.
As you alluded to earlier, we still have to get to orbit.
And then after we get to orbit, we have to really prove out and refine full and rapid reusability.
That ' ll take a moment.
But I do think we will solve this.
I ' m highly confident we will solve this at this point.
CA : Do you ever wake up with \textbf{the} fear that there ' s going to be this Hindenburg moment for SpaceX where...
EM : We ' ve had many Hindenburg.
Well, we ' ve never had Hindenburg moments with people, which is very important.
Big difference.
We ' ve blown up quite a few rockets.
So there ' s a whole compilation online that we put together and others put together, it ' s showing rockets are hard.
I mean, \textbf{the} \textbf{sheer} amount of energy going through a rocket boggles \textbf{the} mind.
So, you know, getting out of Earth ' s gravity well is difficult.
We have a strong gravity and a thick atmosphere.
And Mars, which is less than 40 percent, it ' s like, 37 percent of Earth ' s gravity and has a thin atmosphere. \textbf{The} ship alone can go all \textbf{the} way from \textbf{the} \textbf{surface} of Mars to \textbf{the} \textbf{surface} of Earth.
Whereas getting to Mars requires a giant booster and orbital refilling.
CA : So, Elon, as I think more about this incredible \textbf{array} of things that you ' re involved with, I keep seeing these synergies, to use a horrible word, between them.
You know, for example, \textbf{the} robots you ' re building from Tesla could possibly be pretty handy on Mars, doing some of \textbf{the} dangerous work and so forth.
I mean, maybe there ' s a scenario where your city on Mars doesn ' t need a million people, it needs half a million people and half a million robots.
And that ' s a possibility.
Maybe \textbf{The} Boring Company could play a role helping create some of \textbf{the} subterranean dwelling spaces that you might need.
EM : Yeah.
CA : Back on \textbf{planet} Earth, it seems like a partnership between Boring Company and Tesla could offer an unbelievable deal to a city to say, we will create for you a 3D network of tunnels populated by robo-taxis that will offer fast, low-cost transport to anyone.
You know, full self-driving may or may not be done this year.
And in some cities, like, somewhere like Mumbai, I suspect won ' t be done for a decade.
EM : Some places are more challenging than others.
CA : But today, today, with what you ' ve got, you could put a 3D network of tunnels under there.
EM : Oh, if it ' s just in a tunnel, that ' s a solved problem.
CA : Exactly, full self-driving is a solved problem.
To me, there ' s amazing synergy there.
With Starship, you know, Gwynne Shotwell talked about by 2028 having from city to city, you know, transport on \textbf{planet} Earth.
EM : This is a real possibility. \textbf{The} fastest way to get from one place to another, if it ' s a long distance, is a rocket.
It ' s basically an ICBM.
CA : But it has to land -- Because it ' s an ICBM, it has to land probably offshore, because it ' s loud.
So why not have a tunnel that then connects to \textbf{the} city with Tesla?
And Neuralink.
I mean, if you going to go to Mars having a telepathic connection with loved ones back home, even if there ' s a time delay...
EM : These are not intended to be connected, by \textbf{the} way.
But there certainly could be some synergies, yeah.
CA : Surely there is a growing argument that you should actually put all these things together into one company and just have a company devoted to creating a future that ' s exciting, and let a thousand flowers bloom.
Have you been thinking about that?
EM : I mean, it is tricky because Tesla is a publicly-traded company, and \textbf{the} investor base of Tesla and SpaceX and certainly Boring Company and Neuralink are quite different.
Boring Company and Neuralink are tiny companies.
CA : By comparison.
EM : Yeah, Tesla ' s got 110, 000 people.
SpaceX I think is around 12, 000 people.
Boring Company and Neuralink are both under 200 people.
So they ' re little, tiny companies, but they will probably get bigger in \textbf{the} future.
They will get bigger in \textbf{the} future.
It ' s not that easy to sort of combine these things.
CA : Traditionally, you have said that for SpaceX especially, you wouldn ' t want it public, because public investors wouldn ' t support \textbf{the} craziness of \textbf{the} idea of going to Mars or whatever.
EM : Yeah, making life multi-planetary is outside of \textbf{the} normal time horizon of Wall Street analysts.
To say \textbf{the} least.
CA : I think something ' s changed, though.
What ' s changed is that Tesla is now so powerful and so big and throws off so much cash that you actually could connect \textbf{the} dots here.
Just tell \textbf{the} public that x-billion dollars a year, whatever your number is, will be diverted to \textbf{the} Mars mission.
I suspect you ' d have massive interest in that company.
And it might \textbf{unlock} a lot more possibility for you, no?
EM : I would like to give \textbf{the} public access to ownership of SpaceX, but I mean \textbf{the} thing that like, \textbf{the} overhead associated with a public company is high.
I mean, as a public company, you ' re just constantly sued.
It does occupy like, a fair bit of...
You know, time and effort to deal with these things.
CA : But you would still only have one public company, it would be bigger, and have more things going on.
But instead of being on four boards, you ' d be on one.
EM : I ' m actually not even on \textbf{the} Neuralink or Boring Company boards.
And I don ' t really attend \textbf{the} SpaceX board meetings.
We only have two a year, and I just stop by and chat for an hour. \textbf{The} board overhead for a public company is much higher.
CA : I think some investors probably worry about how your time is being split, and they might be excited by you know, that.
Anyway, I just woke up \textbf{the} other day thinking, just, there are so many ways in which these things connect.
And you know, just \textbf{the} simplicity of that mission, of building a future that is worth getting excited about, might appeal to an awful lot of people.
Elon, you are reported by Forbes and everyone else as now, you know, \textbf{the} world ' s richest person.
EM : That ' s not a sovereign.
CA : EM : You know, I think it ' s fair to say that if somebody is like, \textbf{the} king or de facto king of a country, they ' re wealthier than I am.
CA : But it ' s just harder to measure -- So \$300 billion.
I mean, your net worth on any given day is rising or falling by several billion dollars.
How insane is that?
EM : It ' s bonkers, yeah.
CA : I mean, how do you handle that psychologically?
There aren ' t many people in \textbf{the} world who have to even think about that.
EM : I actually don ' t think about that too much.
But \textbf{the} thing that is actually more difficult and that does make sleeping difficult is that, you know, every good hour or even minute of thinking about Tesla and SpaceX has such a big effect on \textbf{the} company that I really try to work as much as possible, you know, to \textbf{the} edge of sanity, basically.
Because you know, Tesla ' s getting to \textbf{the} point where probably will get to \textbf{the} point later this year, where every high-quality minute of thinking is a million dollars impact on Tesla.
Which is insane.
I mean, \textbf{the} basic, you know, if Tesla is doing, you know, sort of \$2 billion a week, let ' s say, in revenue, it ' s sort of \$300 million a day, seven days a week.
You know, it ' s...
CA : If you can change that by five percent in an hour ' s brainstorm, that ' s a pretty valuable hour.
EM : I mean, there are many instances where a half-hour meeting, I was able to improve \textbf{the} financial outcome of \textbf{the} company by \$100 million in a half-hour meeting.
CA : There are many other people out there who can ' t stand this world of billionaires.
Like, they are hugely offended by \textbf{the} notion that an individual can have \textbf{the} same wealth as, say, a billion or more of \textbf{the} world ' s poorest people.
EM : If they examine sort of -- I think there ' s some axiomatic flaws that are leading them to that conclusion.
For sure, it would be very problematic if I was consuming, you know, billions of dollars a year in personal \textbf{consumption}.
But that is not \textbf{the} case.
In fact, I don ' t even own a home right now.
I ' m literally staying at friends ' places.
If I travel to \textbf{the} Bay Area, which is where most of Tesla engineering is, I basically rotate through friends ' spare bedrooms.
I don ' t have a yacht, I really don ' t take vacations.
It ' s not as though my personal \textbf{consumption} is high.
I mean, \textbf{the} one exception is a plane.
But if I don ' t use \textbf{the} plane, then I have less hours to work.
CA : I mean, I personally think you have shown that you are mostly driven by really quite a deep sense of moral purpose.
Like, your attempts to solve \textbf{the} climate problem have been as powerful as anyone else on \textbf{the} \textbf{planet} that I ' m aware of.
And I actually can ' t understand, personally, I can ' t understand \textbf{the} fact that you get all this criticism from \textbf{the} Left about,"Oh, my God, he ' s so rich, that ' s disgusting."When climate is their issue.
Philanthropy is a topic that some people go to.
Philanthropy is a hard topic.
How do you think about that?
EM : I think if you care about \textbf{the} reality of goodness instead of \textbf{the} perception of it, philanthropy is extremely difficult.
SpaceX, Tesla, Neuralink and \textbf{The} Boring Company are philanthropy.
If you say philanthropy is love of humanity, they are philanthropy.
Tesla is accelerating sustainable energy.
This is a love -- philanthropy.
SpaceX is trying to ensure \textbf{the} long-term survival of humanity with a multiple-planet species.
That is love of humanity.
You know, Neuralink is trying to help solve brain injuries and existential risk with AI.
Love of humanity.
Boring Company is trying to solve traffic, which is hell for most people, and that also is love of humanity.
CA : How upsetting is it to you to hear this constant drumbeat of,"Billionaires, my God, Elon Musk, oh, my God?"Like, do you just shrug that off or does it does it actually hurt?
EM : I mean, at this point, it ' s water off a duck ' s back.
CA : Elon, I ' d like to, as we wrap up now, just pull \textbf{the} camera back and just think...
You ' re a father now of seven surviving kids.
EM : Well, I mean, I ' m trying to set a good example because \textbf{the} birthrate on Earth is so low that we ' re facing civilizational collapse unless \textbf{the} birth rate returns to a sustainable level.
CA : Yeah, you ' ve talked about this a lot, that depopulation is a big problem, and people don ' t understand how big a problem it is.
EM : Population collapse is one of \textbf{the} biggest threats to \textbf{the} future of human civilization.
And that is what is going on right now.
CA : What drives you on a day-to-day basis to do what you do?
EM : I guess, like, I really want to make sure that there is a good future for humanity and that we ' re on a path to understanding \textbf{the} nature of \textbf{the} \textbf{universe}, \textbf{the} meaning of life.
Why are we here, how did we get here?
And in order to understand \textbf{the} nature of \textbf{the} \textbf{universe} and all these fundamental questions, we must expand \textbf{the} scope and scale of consciousness.
Certainly it must not diminish or go out.
Or we certainly won ' t understand this.
I would say I ' ve been motivated by curiosity more than anything, and just desire to think about \textbf{the} future and not be sad, you know?
CA : And are you?
Are you not sad?
EM : I ' m sometimes sad, but mostly I ' m feeling I guess relatively optimistic about \textbf{the} future these days.
There are certainly some big risks that humanity faces.
I think \textbf{the} population collapse is a really big deal, that I wish more people would think about because \textbf{the} birth rate is far below what ' s needed to sustain civilization at its current level.
And there ' s obviously...
We need to take action on climate \textbf{sustainability}, which is being done.
And we need to secure \textbf{the} future of consciousness by being a multi-planet species.
We need to address -- Essentially, it ' s important to take whatever actions we can think of to address \textbf{the} existential risks that affect \textbf{the} future of consciousness.
CA : There ' s a whole generation coming through who seem really sad about \textbf{the} future.
What would you say to them?
EM : Well, I think if you want \textbf{the} future to be good, you must make it so.
Take action to make it good.
And it will be.
CA : Elon, thank you for all this time.
That is a beautiful place to end.
Thanks for all you ' re doing.
EM : You ' re welcome.

% matched lemmas: airplane, array, artificial, asteroid, calculate, carbon, clip, computer, consumption, device, dimension, exact, explosion, fossil, greenhouse, highway, object, planet, probability, sheer, sum, surface, sustainability, terrifying, the, ticket, universe, unlock
\end{textsample}
